% =========================================================
%  INSERT 2 of 3:  Table row ND-13 for tab:future_discriminants
% =========================================================
%  Proposed insertion location:
%  Inside tab:future_discriminants, IMMEDIATELY BEFORE the \bottomrule
%  at line 26021 of ECT_preprint.tex.  The last existing entry is ND12
%  (ending with `\\' on line ~26020).  Insert after that `\\', before
%  \bottomrule on 26021.
%
%  Column structure must match existing entries exactly:
%  # & Observable & ECT statement & Level & Mission/programme
%    & Verdict if data conflict \\
%
%  This row matches that 6-column structure.
% =========================================================

ND13 &
Dark-matter-deficient ultra-diffuse galaxies (DF4, FCC~224; DF2
ambiguous) &
DF4 and probably FCC~224 exhibit tracer kinematics consistent with
baryon-only dynamical mass in the probed region; the current
practical Level-B galactic
closure~(eq.~\ref{eq:ect_rotation_closure_main}) at the baseline
$g^\dagger_{\rm eff}=g^\dagger_0$ tends toward the deep-MOND branch
for these systems and thereby over-predicts the observed
line-of-sight dispersion by a factor of order three.  The minimal
density-only Level-B
ansatz~(eq.~\ref{eq:gdagger_env_main}) for $\Xi$ does not naturally
accommodate both the low-dispersion stress-test objects and matched
DM-rich controls such as NGC~5846-UDG1 within a single simple
environmental scalar law.  The presently known sample is small and
not selection-neutral.  A full diagnostic treatment is given in
Appendix~\ref{app:udg_stress_test}.  Possible theory-level directions
for a future completion (dynamical $G_{\rm eff}(\phi)$ in
eq.~\ref{eq:gen_EE}, history-dependent extensions of $\Xi$, and
orientation-sector coupling) belong to structures already present in
the broader framework, but none is presently derived for the galactic
closure; addressing them is linked to OP3.  $\Lambda$CDM with
bullet-dwarf formation
\cite{Ogiya2018DF2tidal,Moreno2022DMfree,vanDokkum2022bullet}
reproduces the observed dispersions through partial dark-halo
stripping and remains a viable comparator &
B (structural, open) &
Archival: DF4 \cite{Danieli2022DF4},
FCC~224 \cite{Buzzo2025FCC224,Tang2025FCC224},
DF2 \cite{Trujillo2019DF2,Shen2021DF2TRGB},
NGC~5846-UDG1 \cite{Forbes2021NGC5846UDG1};
future: isolated-field DM-deficient searches, IFU $\sigma(R)$
profiles, matched DM-rich controls &
Observation-side: if the sample expands and the proxy-level tension
persists or worsens under same-tracer same-aperture matched
analysis, the minimal Level-B $\Xi$-ansatz is disfavoured and a
completion via $G_{\rm eff}(\phi)$, history-dependent $\Xi$, or the
orientation sector becomes necessary.  Theory-side: a derivation
that maps the dynamical $G_{\rm eff}(\phi)$ degree of freedom onto
the practical galactic closure is required before any ECT rescue
can be claimed \\
